\documentclass[10pt]{article}

\def\Type{LessonCaelan}
\def\TypeNum{Lesson 18}
\def\TypeTitle{Fundamental Theorem of Calculus }
\def\Semester{Fall 2021} %Not Used 
\def\Course{MATH 1190 \& 1210}
\def\Targets{  \bf\color{Fuchsia}I3, \bf\color{Fuchsia}I4 }

\usepackage{preambleDJ}

%\usepackage{ragged2e}


%%%%%%%%%%%% BEGIN DOCUMENT %%%%%%%%%%%%%%%
%%%%%%%%%%%% BEGIN DOCUMENT %%%%%%%%%%%%%%%
%%%%%%%%%%%% BEGIN DOCUMENT %%%%%%%%%%%%%%%
%%%%%%%%%%%% BEGIN DOCUMENT %%%%%%%%%%%%%%%
%%%%%%%%%%%% BEGIN DOCUMENT %%%%%%%%%%%%%%%




\begin{document}

\pagestyle{otherpages}
\thispagestyle{firstpage}
\vs{.65}
Respond to {\bf all 4} numbered prompts to infer a relationship between definite integrals and antiderivatives!

\begin{enumerate}
\item Consider the function $f(x) = x-2$. 
    \begin{enumerate}
    \item The graph of $f$ is given below. Calculate $\dint_{0}^{4}f(x)\, dx$ by interpreting the definite integral as net area. (See lesson 16 if you need a refresher on this!) Show your work. \\
    \mpageT{.35}{.65}{
        \begin{tikzpicture}[scale=1]
              \begin{axis}[
                xlabel={$x$},
                %grid = major,
                xtick       = {1,2,3,4,5},
                xticklabels = { ,$2$, ,$4$},
                ytick       = {-2,-1,1,2,3,4,5},
                yticklabels = {$-2$, , ,$2$, ,$4$, },
                samples     = 320,
                domain      = -1:9,
                restrict y to domain=-40:40,
                xmin = -1, xmax = 5,
                ymin = -3, ymax = 6,
                height=2.5in, width=2.5in,
                grid=both,grid style=dashed
              ]
              % old version of the curve
              %
              %\addplot[domain=-1:9, purple, ultra thick, mark=none] {-1/10*(x-1)*(x-4)^2*(x-7)};
              %
              %new version of the curve in pieces
              %
              \addplot[domain=-1:5, ultra thick, red, ] {x-2};
              \addplot[name path=f1, domain=0:2, ultra thick, red, ] {x-2};
  \path[name path=axis1]  (axis cs:0,0)--(axis cs:2,0) ;
     \addplot [
        thick,
        color=blue,
        fill=blue, 
        fill opacity=0.1
    ]
    fill between[
        of=f1 and axis1,
        soft clip={domain=0.:2},
   ];
   \addplot[name path=f2, domain=2:4, ultra thick, red, ] {x-2};
 \path[name path=axis2]  (axis cs:2,0)--(axis cs:4,0) ;
     \addplot [
        thick,
        color=red,
        fill=red, 
        fill opacity=0.1
    ]
    fill between[
        of=f2 and axis2,
        soft clip={domain=2:4},
   ];
            \end{axis}
        \end{tikzpicture}
}{
\sol{ Since part of the area is below the $x$-axis and some above, we use integral properties to split up the integral and use area of triangles to find net area.
\itemI{
\item $\ds \int_0^4 f(x)\;dx=\int_0^2f(x)\;dx+\int_2^4 f(x)\;dx$
\item $\ds \int_0^2f(x)\;dx= -\frac{1}{2}b\cdot h=-\frac{1}{2}2\cdot 2=-2$. Note answer is negative because area is below $x$-axis.  
\item  $\ds \int_2^4f(x)\;dx= \frac{1}{2}b\cdot h=-\frac{1}{2}2\cdot 2=2$
\item $\ds \int_0^4 f(x)\;dx=\int_0^2f(x)\;dx+\int_2^4 f(x)\;dx= (-2)+2=0$
}
}
}    
  \item Find an antiderivative of $f$: \quad \dotfill \quad $F(x)=$\unline{2}{\blue{$\ds \frac{x^2}{2}-2x$}}{.25}

    \item Compute the net change in your antiderivative as $x$ changes from $0$ to $4$. That is, find $\Delta F = F(4) - F(0)$. Show your work. 
\sol{
\al{
\Delta F&=F(4)-F(0)\\
&=\left(\frac{4^2}{2}-2\cdot 4 \right)-\left( \frac{0^2}{2}-2\cdot0\right)\\
&=(8-8)-(0-0)\\
&=0
}
}  
    \item Compare your answers in (a) and (c). Specifically, how are $F(4)-F(0)$ and $\dint_{0}^{4}f(x)\,dx$ related? 
       \sol{
       $$\ds \int_{0}^{4}f(x)\,dx=F(4)-F(0)$$ 
       }
    \item Describe in words how you can use an antiderivative of a function to find a definite integral of that function if the relationship you described in part (d) is true more generally.
\sol{ Let's assume our definite integral is of the form $\ds \int_a^b f(x)\;dx$.  Steps should include
\itemI{
\item Find an anti-derivative of the function $f(x)$. We will call it $F(x)$
\item Substitute the upper and lower bounds into the anti-derivative $F(x)$.
\item Subtract the two values (value from upper bound - value from lower bound, i.e. $F(b)-F(a)$ )
}
Note that using correct words alone without math symbols can receive full credit if explanation is equivalent.
}
        
    \end{enumerate}

%%%%%%%%%
\newpage%
%%%%%%%%%

\item Find $\dint_0^4 3 \,dx$ using the method you discovered in prompt \#1. Show your work.
\sol{
 \itemI{
 \item $\ds \int_a^bf(x)=\;dx =\int_0^4 3 \,dx$. Thus $f(x)=3$.
 \item An antiderivative of $f(x)=3$ is $F(x)=3x$.
 \item Substitute upper and lower bounds.  $F(4)=3\cdot 4$.  $F(0)=3\cdot 0$.
 \item Subtract: $$\int_0^4 3 \,dx=F(4)-F(0)=12-0=12$$
 }
}
\item Find $\dint_0^2 (4-2x) \,dx$ using the method you discovered in prompt \#1. Show your work.
\sol{
 \itemI{
 \item $\ds \int_a^bf(x)=\;dx =\int_0^2 4-2x \,dx$. Thus $f(x)=4-2x$.
 \item An antiderivative of $f(x)=4-2x$ is $F(x)=4x-x^2$.
 \item Substitute upper and lower bounds.  $F(2)=4\cdot2-2^2=8-4=4$.  $F(0)=4\cdot0-0^2==$.
 \item Subtract: $$\int_0^3 4-2x \,dx=F(2)-F(0)=4-0=4$$
 }
}

\item If $F$ is an antiderivative of $f$, then what expression involving $F$ completes the formula below?

\[ \dint_a^b f(x)\,dx =\unline{2}{\blue{$F(b)-F(a)$}}{.25}  \]

\

\


\end{enumerate}

\end{document}\end{document}
